\chapter{Сигнал тривоги}
% full version: chapters/14_pain.tex

\pencilfig{14}{\pencilart{14}}{Біль --- пожежна сигналізація, а не термометр.}

Усі інші чуття повідомляють, яким є світ. Біль --- ні. Він повідомляє, що щось загрожує тілу, і його гучність багато в чому задає сама машина, а не подразник. Біль --- це пожежна сигналізація, а не термометр.

Датчики болю спрацьовують лише на сильне: на жар, вищий за сорок із лишком градусів, на сильний тиск, на їдкі речовини. Вони звикають значно слабше, ніж датчики дотику, --- сигналізація не повинна замовкати, поки загроза є.

\section*{Два проводи}

Згадайте, як ви вдарили палець. Спочатку --- гострий, точний біль: де і коли. За мить --- тупий, розлитий, тривалий: наскільки погано. Це два різні проводи. Один швидкий, його сигнал долітає за десятки мілісекунд. Інший повільний, його сигнал іде близько секунди.

\section*{Ворота}

Чому, забившись, ми тремо забите місце? У спинному мозку сигнал болю проходить через «ворота». Дотик будить гальмівний нейрон, який їх прикриває. Тому потерти забите місце справді допомагає --- але лише проти помірного болю; сильний однаково проходить. Що сильний біль при цьому сам вимикає гальмівний нейрон, --- частина вихідної теорії воріт, її прямого підтвердження поки немає.

\pencilfig{126}{%
% gate: the pain wire drives the relay neuron; touch wakes an inhibitory neuron that closes the gate
\draw[pen=pcAmber] (0,1.35) -- (3.05,1.0); \draw[arr=pcAmber] (2.8,1.03) -- (3.08,1.0);
\node[lbl, text=pcAmber!70!black, anchor=east] at (-0.05,1.35) {біль};
\draw[pen=pcBlue] (0,-0.15) -- (1.75,-0.15); \draw[arr=pcBlue] (1.5,-0.15) -- (1.78,-0.15);
\node[lbl, text=pcBlue, anchor=east] at (-0.05,-0.15) {дотик};
\draw[dotpen=pcRed, wash=pcRed] (2.05,-0.15) circle (0.26);
\draw[pcRed, line width=0.7pt, -{Bar[width=6pt]}] (2.25,0.05) -- (3.2,0.66);
\node[lbl, text=pcRed, anchor=north] at (2.05,-0.45) {гальмо};
\draw[dotpen=pcGray, wash=pcGray] (3.4,0.9) circle (0.34);
\node[lbl, anchor=south] at (3.4,1.3) {ворота};
\draw[pen] (3.75,0.9) -- (5.6,0.9); \draw[arr] (5.35,0.9) -- (5.65,0.9);
\node[lbl, anchor=west] at (5.7,0.9) {у мозок};
}{Ворота в спинному мозку. Сигнал болю проходить у мозок через нейрон-ворота, а дотик будить гальмівний нейрон, який їх прикриває. Тому потерти забите місце допомагає, але лише проти помірного болю.}

\section*{Ручка гучності згори}

Мозок уміє сам робити біль гучнішим чи тихішим. Згори, з головного мозку, до воріт ідуть радіоканали --- \nt{SERO}, \nt{NORA} та особливі речовини, схожі на знеболювальні, які мозок виділяє сам. У небезпечній ситуації біль може майже зникнути: тому, хто тікає від небезпеки, ніколи кульгати. Очікування полегшення теж приглушує біль, почасти через ті самі речовини.

\section*{Коли сигналізація залипає}

Повторний біль підвищує чутливість: той самий удар починає боліти сильніше. Це теж навчання --- петля, яка сама себе підкручує, поки рана гоїться. Але сильна петля має небезпеку: вона може «залипнути» й далі дзвонити, коли рана вже загоїлася, --- як група нейронів із другого розділу, що підтримує сама себе.

\section*{Учитель і переривання}

Біль --- ще й учитель: для дофамінового «краще чи гірше, ніж очікувалося» він працює як від'ємна винагорода. І він перериває поточну справу, як спалах норадреналіну: кидай усе, займися небезпекою. А найшвидша реакція --- відсмикнути руку від гарячого --- замикається просто в спинному мозку тією самою парою м'язів-супротивників із попереднього розділу, ще до того, як ви встигли відчути біль.

\fullversion{14}
